\section{Background and Related Work}
\label{sec:related_work}

Though hallucination is often deemed to be essential for the creativity of LMs \citep{math11102320, 10.1145/3652102}, it remains one of the most frequently criticized shortcomings of these models. While some studies argue that hallucinations are an inherent consequence of their stochastic nature, often known as the `stochastic parrots' phenomenon \citep{10.1145/3442188.3445922}, many others attribute the issue primarily to artifacts in the data that these models ingest. There is a substantial body of research addressing diverse aspects of hallucination in LMs -- we briefly summarize a few of these works related to two major dimensions: sources of hallucination and hallucination detection.

\subsection{Sources of Hallucination} 
% The sources of hallucination can be broadly classified into three categories: \textit{Data, Training and Inference}~\citep{Huang_2025}.

% \paragraph{Data-related hallucination} Although pre-training is a crucial phase in the creation of LLMs, it may also be a cause of hallucination, especially owing to flawed data sources and inferior data utilization. With the scaling of pre-training data with ever larger models, misinformation and biases within the data leads to imitative falsehoods~\citep{lin-etal-2022-truthfulqa}, duplication bias shifting the LLM from generalization to memorization~\citep{lee_deduplicating_2022, hernandez2022scalinglawsinterpretabilitylearning}, and social biases such as related to gender~\citep{data_paullada_2021} and nationality~\citep{narayanan-venkit-etal-2023-nationality, ladhak-etal-2023-pre}. With time, LLMs may also start lagging behind factual knowledge , which becomes more substantial while addressing domain-specific queries~\citep{kasai2024realtimeqawhatsanswer}.

% \paragraph{Training-related Hallucination} At pre-training stage, design choices like inherent architectural design may cause flaws like inadequate unidirectional representation~\citep{li2023batgptbidirectionalautoregessivetalker} and attention glitches~\citep{liu2023exposingattentionglitchesflipflop}, along with exposure bias~\citep{wang_exposure_2020} pose a hallucination risk. Meanwhile at post-training stage - alignment of LLMs may cause hallucination related to misalignment of capability and belief~\citep{sharma_towards_2023}.

% \paragraph{Inference-related Hallucination} Decoding algorithms such as Stochastic sampling~\citep{fan-etal-2018-hierarchical, holtzman_curious_2019} introduce unpredictability in LLM creation, resulting in creativity and diversity. However, this incurs a heightened risk of hallucination~\citep{dziri-etal-2021-neural, chuang2024doladecodingcontrastinglayers}. Other flaws like insufficient context attention~\citep{liu_instruction_2023, shi2023trustingevidencehallucinatecontextaware} and softmax bottleneck~\citep{yang2018breakingsoftmaxbottleneckhighrank, chang_softmax_2022} in the top-layer representations also pose a challenge. 

Hallucinations in LMs have been attributed to a complex interplay of factors spanning data quality, model architecture, and inference processes. On the data front, pre-training on vast and often imperfect corpora introduces risks of propagating misinformation, factual inaccuracies, and social biases~\citep{lin-etal-2022-truthfulqa, lee_deduplicating_2022, hernandez2022scalinglawsinterpretabilitylearning, data_paullada_2021, narayanan-venkit-etal-2023-nationality, ladhak-etal-2023-pre, kasai2024realtimeqawhatsanswer}. The ever-increasing scale of training data, while enabling greater generalization, can also amplify errors and result in LMs memorizing, rather than reasoning about information.

At the level of model training and design, certain architectural limitations -- such as insufficient bidirectional representation or attention mechanism failures -- can introduce systematic vulnerabilities to hallucination~\citep{li2023batgptbidirectionalautoregessivetalker, liu2023exposingattentionglitchesflipflop, wang_exposure_2020, sharma_towards_2023}. Moreover, the alignment process that tunes LMs to follow human preferences may occasionally introduce a mismatch between model capabilities and desired behaviors, leading to unexpected model outputs.

Inference strategies, especially those introducing stochasticity or prioritizing diversity (such as top-$k$ or nucleus sampling), have also been linked to hallucination~\citep{fan-etal-2018-hierarchical, holtzman_curious_2019, dziri-etal-2021-neural, chuang2024doladecodingcontrastinglayers}. Additional challenges arise from limitations in context integration and representational bottlenecks in output layers~\citep{liu_instruction_2023, shi2023trustingevidencehallucinatecontextaware, yang2018breakingsoftmaxbottleneckhighrank, chang_softmax_2022}. Together, these findings indicate that hallucinations arise not from a single source but from interactions among data, architecture, and inference.


\subsection{Hallucination Detection} 
% Hallucinations demonstrated by LLMs have attracted significant interest because of their impact on model dependability and practical use. As models improve in producing human-like writing, differentiating between factual and fabricated information becomes a critical issue. Due to broad nature of hallucinations, various approaches can be used. Factuality Hallucination can be countered with techniques like \textit{external knowledge verification}, which checks outputs against external sources \citep{chen2024complexclaimverificationevidence, min2023factscorefinegrainedatomicevaluation, huo2023retrievingsupportingevidencellms, chern2023factoolfactualitydetectiongenerative} but faces resource maintenance challenges. \textit{Faithfulness evaluation} techniques use fact-based overlap metrics \citep{nan-etal-2021-entity}, NLI-based classifiers \citep{falke-etal-2019-ranking, maynez-etal-2020-faithfulness}, or QA-based consistency checks \citep{durmus-etal-2020-feqa, scialom-etal-2021-questeval}. Another branch of research argues that origin of LLM hallucination is correlated with model's uncertainty, as a result tries to find way of estimating uncertainty. These can be further classified depending on availability of LLM's internal states. If LLM's internal states can be accessed, uncertainty can be measured using entropy-based methods~\citep{xiao-wang-2021-hallucination, malinin_uncertainty_2021} and log-probability based~\citep{guerreiro-etal-2023-looking}. Some LLMs are black box and accessible only through API calls. For such cases, uncertainty can be estimated through natural language prompts\citep{xiong2024llmsexpressuncertaintyempirical, kadavath2022languagemodelsmostlyknow} or techniques like SelfCheckGPT which samples multiple generations from an LLM for the same prompt\citep{manakul2023selfcheckgptzeroresourceblackboxhallucination}. Research also proposes using LLMs's capabilities as evaluators to asses the model-generated content~\citep{chiang2023largelanguagemodelsalternative, luo2023chatgptfactualinconsistencyevaluator, laban2023llmsfactualreasonersinsights, miao2023selfcheckusingllmszeroshot}

Given the broad impact of hallucination, a wide range of detection strategies has been proposed. Many early efforts focused on evaluating model confidence, with the intuition that LMs are less certain when hallucinating. Metrics such as perplexity, entropy, and energy-based scores~\citep{ren_out--distribution_2023, malinin_uncertainty_2021, liu_energy-based_2020} have been explored for this purpose. Another family of approaches leverages output consistency: by sampling multiple generations for the same prompt, one can measure lexical or semantic agreement across responses~\citep{lin-etal-2022-truthfulqa, durmus-etal-2020-feqa, scialom-etal-2021-questeval}. High variability is often taken as a proxy for hallucination, and techniques such as BERTScore, NLI-based classification, and question-answering-based evaluation have all been applied~\citep{falke-etal-2019-ranking, maynez-etal-2020-faithfulness, nan-etal-2021-entity, chen2024complexclaimverificationevidence}. However, these multi-pass strategies are computationally demanding and poorly suited to low-latency settings, as they require generating and analyzing numerous outputs for each input.

Other detection frameworks rely on verifying factuality by grounding model outputs in external knowledge sources~\citep{min2023factscorefinegrainedatomicevaluation, huo2023retrievingsupportingevidencellms, chern2023factoolfactualitydetectiongenerative, chen2024complexclaimverificationevidence}. While these methods can be effective, they require access to curated external databases and ongoing maintenance, and may not generalize across domains.

More recently, studies have explored the use of LMs themselves as evaluators -- either via prompt engineering that elicits self-critique, or by prompting additional LMs to judge the factuality or faithfulness of generated content~\citep{xiong2024llmsexpressuncertaintyempirical, kadavath2022languagemodelsmostlyknow, manakul2023selfcheckgptzeroresourceblackboxhallucination, chiang2023largelanguagemodelsalternative, luo2023chatgptfactualinconsistencyevaluator, laban2023llmsfactualreasonersinsights, miao2023selfcheckusingllmszeroshot}. These meta-evaluation techniques provide flexibility in black-box settings, but can inherit the blind spots of the underlying models and are themselves subject to hallucination.

% A comparatively smaller but growing line of work examines the internal dynamics of language models, investigating whether the hidden states or intermediate representations carry signals that are indicative of hallucination~\citep{chen2024complexclaimverificationevidence, ji-etal-2024-llm}. These studies have shown that certain patterns of representational diversity of the hidden states may be predictive of hallucinated content. However, most prior work in this area has either required separate classifier training on internal states or operated in multi-pass setups, limiting their practicality for inference-time detection.

% \changed{
% \paragraph{Mechanistic Interpretability Methods}
% A growing body of work investigates hallucinations by training classifiers (``probes'') on internal representations to predict factuality or truthfulness ~\citep{azaria2023internalstatellmknows, ji-etal-2024-llm, obeso2025realtimedetectionhallucinatedentities}. These methods demonstrate that hidden states encode useful signals for hallucination detection, but typically require supervised training on task- and model-specific data. As a result, their effectiveness can degrade under domain or task shifts, and their deployment requires labeled data and retraining for new settings.}
% \changed{
% In contrast, \method~does not train any auxiliary model and instead directly measures statistical dependence between input and output representations at inference time. This distinction is crucial: \method~is designed to be training-free and task-agnostic, enabling deployment without access to labeled hallucination data. In Section~\ref{sec:results}, we empirically compare \method~against a strong supervised probing baseline and show that, while probes perform well on their training domain, \method~exhibits more robust generalization to unseen datasets without any training.
% }

\changed{
\paragraph{Mechanistic Interpretability and Representation-based Hallucination Detection}
A growing body of work in mechanistic interpretability investigates how hallucinations are reflected in the internal representations of language models. Several studies demonstrate that hidden states encode signals predictive of factuality or truthfulness, and leverage supervised probes to extract these signals for hallucination detection or analysis \citep{azaria2023internalstatellmknows, ji-etal-2024-llm, duan2024llmsknowhallucinationempirical, obeso2025realtimedetectionhallucinatedentities}. Recent work has also explored latent-space diversity, activation geometry, and representation dynamics as indicators of hallucinated generation \citep{chen_inside:_2024, geva-etal-2023-dissecting}. Collectively, these studies provide strong evidence that hallucinations are systematically encoded within LM hidden states.
}
\changed{
While effective, most representation-based detection methods rely on supervised training of task- and model-specific classifiers, or require dataset-dependent calibration. As a result, their performance can degrade under domain or task shifts, and their deployment typically requires access to labeled hallucination data and retraining for new settings. In contrast, \method~does not train any auxiliary model and instead directly measures statistical dependence between input and output representations at inference time. This design makes \method~training-free and task-agnostic, enabling deployment without labeled hallucination data. In Section~\ref{sec:results}, we empirically compare \method~against a strong supervised probing baseline and show that, while probes perform well on their training domain, \method~exhibits more robust generalization to unseen datasets without any training.
}


\subsection{Background of Hallucination Detection Measures} \label{sec:background}

Here we briefly discuss the mathematical formulation of a few hallucination detection measures, which are used as baselines in this work. Since our proposed method doesn't involve any training, we compare only against other training-free methods.


Assume that $\mathbf{S}_{in} = \{s_1, s_2, \ldots, s_I\}$ is the input prompt, and $\mathbf{S}_{out} = \{t_1, t_2, \ldots, t_O\}$ is the corresponding output generated by the LM, parameterized by $\theta$, where $\mathbf{S}_{out}$ is a sequence of $O$ tokens. In case of measures requiring multiple generations for the same input $\mathbf{S}_{in} $, assume $\mathcal{S} =[\mathbf{S}_{out}^1, \mathbf{S}_{out}^2, ..., \mathbf{S}_{out}^N]$ to be the collection of $N$ generations, where diversity is ensured through stochastic sampling strategies, like top-$k$, nucleus or temperature sampling. Also, $\mathbb{P}_{LM}(.)$ denotes the probability distribution generated by the LM for any position in the sequence. 
\paragraph{\textbf{Perplexity}}
 Perplexity, as a token-level uncertainty measure, can be utilized for evaluating potential hallucinations \citep{ren_out--distribution_2023}. Perplexity is monotonically related to the mean negative log-likelihood (MNLL) of the output sequence, and hence MNLL can be utilized as a tractable proxy for perplexity as a hallucination detection measure \citep{chen_inside:_2024}. Formally, MNLL can be defined as: $\text{MNLL}(\mathbf{S}_{out}|\mathbf{S}_{in}; \theta) = -\frac{1}{O}\sum_{i=1} ^{O}\log \mathbb{P}_{LM}(t_i|t_{<i}, \mathbf{S}_{in})$. A lower value of this measure indicates that the model is more confident about its generation. %While perplexity is simple to implement, it has notable limitations. Since shorter sequences generally have lower perplexity, the length of the output sequence $T$ is used to normalize the joint probability. More significantly, perplexity defined by averaging token-level uncertainty cannot effectively capture the uncertainty of the entire sequence, as different tokens contribute differently to the semantics of the sentence \citep{duan2024shiftingattentionrelevancepredictive, raj2025semanticconsistencyassuringreliability}.

\paragraph{\textbf{Energy Score}}
The Energy Score \citep{liu_energy-based_2020} approaches hallucination detection from a different perspective by leveraging energy-based models. Unlike softmax-based confidence scores, energy scores are aligned with the probability density of the inputs. 
%The energy function maps each input to a single scalar that is lower for observed data and higher for unobserved ones. 
Let $\mathbb{L}_{LM}(.)$ denote the logit distribution over the set of tokens in the vocabulary $\mathcal{V}$, generated by the LM, such that $\mathbb{P}_{LM}(.) = \text{softmax}(\mathbb{L}_{LM}(.))$. The energy score, for a temperature parameter $T$, can then be defined as: $\text{Energy}(\mathbf{S}_{in}; \theta) = -T \cdot \log\sum_{t \in \mathcal{V}} e^{\mathbb{L}_{LM}(t|\mathbf{S}_{in})/T}$. The energy score is generally used for out-of-distribution detection, where samples with higher energies can be interpreted as datapoints with a lower likelihood of occurrence.

\paragraph{\textbf{Length-Normalized Entropy}} Unlike perplexity which utilizes a single generation for quantifying sequence-level uncertainty, length-normalized entropy (LN-Entropy) uses multiple generations to measure the degree of hallucination~\citep{malinin_uncertainty_2021}. LN-Entropy is defined as: $\text{LNE}(\mathcal{S}|\mathbf{S}_{in}; \theta) = -\mathbb{E}_{\mathbf{S}_{out}\in \mathcal{S}}\frac{1}{|\mathbf{S}_{out}|}\sum_{i=1}^{|\mathbf{S}_{out}|}\log \mathbb{P}_{LM}(t_i|t_{<i}, \mathbf{S}_{in})$. It is assumed that when an LM generates hallucinated content, it is typically uncertain, resulting in an output distribution with higher entropy \cite{kadavath2022languagemodelsmostlyknow}.

\paragraph{\textbf{Lexical Similarity}} Lexical Similarity (LS) \cite{lin_towards_2022} aims to quantify the similarities between the multiple generations from an LM corresponding to the same input prompt, with the assumption that LS will be low in case of hallucinations. It can formally defined as: 
$\text{LS}(\mathcal{S}|\mathbf{S}_{in}; \theta) = \frac{2}{N(N-1)}\sum_{i=1}^{N}\sum_{j=i+1}^{N}sim(\mathbf{S}_{out}^i, \mathbf{S}_{out}^j)$, where, $sim(\cdot, \cdot)$ can be any function to quantify similarity, like Rouge-L \cite{lin-2004-rouge}.

\paragraph{\textbf{EigenScore}}
EigenScore \citep{chen_inside:_2024} leverages the hidden states of LMs for hallucination detection. It measures semantic diversity in the latent space of LMs' internal representations by calculating the log-determinant of the covariance matrix of multiple output representations: $\text{ES}(\mathcal{S}|\mathbf{S}_{in}; \theta)  = \frac{1}{N}\log\det(\Sigma + \alpha \cdot \mathbb{I}_N) = \frac{1}{N}\sum_{i=1}^{N}\log(\lambda_i)$, where, the covariance matrix of the output representations is given by: $\Sigma = Z^T \cdot C_d \cdot Z$ with $Z = [z_1, z_2, \cdots, z_N] \in \mathbb{R}^{d\times N}$ being the embedding matrix of the $N$ generated outputs and $C_d = \mathbb{I}_d - \frac{1}{d}\mathbf{1}_d\mathbf{1}_d^T$ being the centering matrix. The set of eigenvalues of the regularized covariance matrix $\Sigma + \alpha \cdot \mathbb{I}$ is given by $\lambda = \{\lambda_1, \lambda_2, \cdots, \lambda_N\}$. It is assumed that when the LM is not hallucinating, the representations of the outputs will be highly correlated, thereby having $\lambda_i=0$ for most $i$'s.

Table~\ref{tab:evaluation_methods} presents a summary of the characteristics of the discussed hallucination detection measures, and compares our proposed method with these techniques.

\begin{table}[t!]
\centering
\resizebox{0.8\linewidth}{!}{
\begin{tabular}{lccc}
\toprule
\textbf{Method} & \makecell{\textbf{Uses} \\ \textbf{Uncertainty}} & \makecell{\textbf{Checks} \\ \textbf{Consistency}} & \makecell{\textbf{No. of Generations} \\ \textbf{Required}} \\
\midrule
Perplexity       & \cmark & \xmark & 1 \\
Energy        & \cmark & \xmark & 1 \\
Length-Normalized Entropy     & \cmark & \xmark & $N$ \\
Lexical Similarity & \xmark & \cmark & $N$ \\
Eigenscore & \xmark & \cmark & $N$ \\
\rowcolor{gray!10}
\textbf{\method~(Ours)}        & \xmark & \xmark & \text{1} \\
\bottomrule
\end{tabular}
}
\caption{Comparison of some popular training-free hallucination detection techniques. Most methods rely on either model uncertainty or consistency across multiple generations ($N$) -- approaches relying on the latter show increased latency and compute cost. \method, on the other hand, requires only a single generation and leverages internal representation independence, offering a lightweight yet effective alternative.}
\label{tab:evaluation_methods}
\end{table}